\documentclass{article}
\usepackage[utf8]{inputenc}
\usepackage{enumitem}
\title{Forslag til eksperimenter}
\author{}
\date{}

\begin{document}
\maketitle

\begin{enumerate}
\item \textbf{Måle toppstrøm:} Bruk en Rogowski-spole eller en strømtransformator rundt en av ledningene til spolen. Integrer signalet for å finne strømmen. Plott strøm som funksjon av tid. Forventet toppstrøm: 8--9 kA.
\item \textbf{Kartlegge magnetfelt:} Bruk en liten pickup-spole (1 cm$^2$) og mål indusert spenning som funksjon av avstand og vinkel. Plott $V_{\text{ind}}$ vs. $1/r^3$ for å verifisere dipolmodellen. Sammenlign med teoretisk feltstyrke (0.68 T i sentrum).
\item \textbf{Sammenligne spoledesign:} Test ulike antall vindinger (4, 6, 8) og diametre. Mål resonansfrekvens og toppstrøm. Finn optimal konfigurasjon for maksimalt magnetfelt.
\item \textbf{Teste skjerming:} Plasser forskjellige materialer (aluminiumsfolie, kobbernett, stålplate) mellom spolen og en mottakerkrets. Mål dempingen av den induserte spenningen.
\item \textbf{Ødeleggelsesterskler:} Finn minimum avstand for å ødelegge ulike enheter (LED, transistor, mikrokontroller). Dokumenter resultatene.
\item \textbf{Repetisjonsrate:} Mål hvor raskt du kan fyre med kontinuerlig lading. Varier batterispenning og ladestrøm for å optimalisere.
\item \textbf{Stigetidsmåling:} Bruk en høyhastighets pickup-spole og oscilloskop for å måle den faktiske stigetiden til magnetfeltet. Sammenlign med teoretisk verdi (~70 µs).
\end{enumerate}
\end{document}
